\documentclass[10pt]{article}
\usepackage[T1]{fontenc}
\usepackage[utf8]{inputenc}
\usepackage[french]{babel}
\usepackage{amsmath,amssymb,mathrsfs,mathpazo}
\usepackage{geometry}
\geometry{a4paper,margin=2cm}
\usepackage{graphicx,array,tabularx,tabularray,colortbl,multirow,multicol,booktabs,enumitem,pifont,ulem,fancybox,diagbox,cancel}
\usepackage{pstricks,pstricks-add,pst-plot,pst-tree,pst-node,pst-text,pst-func,pst-math,pst-bspline,pst-3dplot,pst-eucl,pst-solides3d}
\let\pscalculate\undefined
\usepackage{pst-bezier}
\usepackage{tikz,pgfplots,tkz-tab}
\usetikzlibrary{arrows,arrows.meta,patterns,calc,positioning,trees,decorations.markings,shapes.geometric,shapes.misc,shapes.arrows,intersections,decorations.pathmorphing,decorations.pathreplacing,backgrounds,shadows,babel}
\pgfplotsset{compat=1.18}
\usepgfplotslibrary{fillbetween}
\usepackage[np]{numprint}
\usepackage{listings,makecell,pas-tableur,twoopt}
\let\barre\undefined
\usepackage[pstricks]{tablvar}
\definecolor{darkgreen}{RGB}{0,100,0}
\providecommand{\red}{\color{red}}
\providecommand{\blue}{\color{blue}}
\providecommand{\green}{\color{green}}
\providecommand{\black}{\color{black}}
\providecommand{\white}{\color{white}}
\providecommand{\gray}{\color{gray}}
\providecommand{\cyan}{\color{cyan}}
\providecommand{\magenta}{\color{magenta}}
\providecommand{\yellow}{\color{yellow}}
\providecommand{\violet}{\color{violet}}
\providecommand{\orange}{\color{orange}}
\providecommand{\euro}{\texteuro}
\providecommand{\EUR}{\texteuro}
\providecommand{\ieme}{\textsuperscript{e}}
\providecommand{\degres}{\ensuremath{^\circ}}
\providecommand{\degree}{\ensuremath{^\circ}}
\providecommand{\vv}[1]{\overrightarrow{#1}}
\providecommand{\Coord}[1]{#1}
\providecommand{\decosix}{}
\providecommand{\starredbullet}{\ensuremath{\bullet}}
\pagestyle{empty}
\setlength{\parindent}{0pt}
\usepackage{adjustbox}\DeclareUnicodeCharacter{2236}{:}\DeclareUnicodeCharacter{2190}{\ensuremath{\leftarrow}}
\providecommand{\R}{\mathbb{R}}
\providecommand{\N}{\mathbb{N}}
\providecommand{\Z}{\mathbb{Z}}
\providecommand{\Q}{\mathbb{Q}}
\providecommand{\D}{\mathbb{D}}
\providecommand{\C}{\mathbb{C}}
\providecommand{\e}{{\mathrm{e}}}
\providecommand{\vect}[1]{\overrightarrow{#1}}
\providecommand{\vectt}[1]{\overrightarrow{\mathrm{#1}}}
\providecommand{\vvt}[1]{\overrightarrow{\mathrm{#1}}}
\providecommand{\np}[2][]{{#2~#1}}
\providecommand{\vv}[1]{{\overrightarrow{#1}}}
\renewcommand{\d}{{\mathrm{d}}}
\providecommand{\degres}{{^\circ}}
\providecommand{\degre}{{^\circ}}
\providecommand{\blue}{\color{blue}}
\providecommand{\red}{\color{red}}
\providecommand{\enskip}{\quad}
\providecommand{\llbracket}{\lbrack\!\lbrack}
\providecommand{\rrbracket}{\rbrack\!\rbrack}
\providecommand{\euro}{{\text{€}}}
\providecommand{\no}{{\text{n°}}}
\providecommand{\interval}[2]{{[#1;#2]}}
\providecommand{\intervFO}[2]{{[#1;#2[}}
\providecommand{\dotfill}{\ldots}

\providecommand{\ds}{\displaystyle}
\providecommand{\textding}[1]{#1}
\providecommand{\Oij}{\(\mathrm{O};\vec{\imath},\vec{\jmath}\)}
\providecommand{\Oijk}{\(\mathrm{O};\vec{\imath},\vec{\jmath},\vec{k}\)}
\providecommand{\Ouv}{\(\mathrm{O};\vec{u},\vec{v}\)}
\providecommand{\intervFF}[2]{\ensuremath{[#1;#2]}}
\renewcommand{\euro}{\ifmmode\text{\texteuro}\else\texteuro\fi}
\DeclareUnicodeCharacter{00B0}{\ifmmode^{\circ}\else\textdegree\fi}
\DeclareUnicodeCharacter{20AC}{\ifmmode\text{\texteuro}\else\texteuro\fi}
% Énoncés : Éducation nationale. Transcription et corrigés : APMEP, Denis Vergès et contributeurs. Découpage et classement : CoopMaths, Sébastien Lozano, Éric Elter et contributeurs. Adaptation de la présentation : Maths Explorer.
% CC BY-SA : https://coopmaths.fr/www/about/
% https://coopmaths.fr/alea/static/bac/2023/tex/bac_2023_03_sujet2_ameriquenord_1_cor.tex
\begin{document}
\noindent\textit{APMEP / CoopMaths. CC BY-SA. Présentation adaptée par Maths Explorer.}\par\medskip
\medskip

\textbf{\large Partie A}

\medskip

Le plan est muni d'un repère orthogonal.

On considère une fonction $f$ définie et dérivable sur $\mathbb{R}$. On note $f'$ sa fonction dérivée. On donne ci-dessous la courbe représentative de la fonction dérivée $f'$.

\begin{center}
\scalebox{0.8}{
\psset{unit=1cm,arrowsize=2pt 3}
\begin{pspicture*}(-7,-8)(3,3)
\psaxes[linewidth=1.25pt,labelFontSize=\scriptstyle]{->}(0,0)(-7,-7.95)(3,3)
\psgrid[gridlabels=0pt,subgriddiv=1,gridwidth=0.3pt,gridcolor=gray]
\psplot[plotpoints=2000,linewidth=1.25pt,linecolor=red]{-7}{3}{x dup mul 3 x mul sub 1 add 2.71828 x exp mul}
\end{pspicture*}
}
\end{center}





\begin{enumerate}
\item Le sens de variation de la fonction $f$ sur $\mathbb{R}$ est donné par le signe de la dérivée $f'$.

~~\textbullet~~Sur $]-\infty~;~0,4[$, $f'>0$ donc $f$ est strictement croissante.

~~\textbullet~~Sur $]0,4~;~2,6[$, $f'<0$ donc $f$ est strictement décroissante.

~~\textbullet~~Sur $]2,6~;~+\infty[$, $f'>0$ donc $f$ est strictement croissante.

\item La fonction $f$ est convexe sur les intervalles sur lesquels la dérivée $f'$ est croissante, soit sur $]-\infty~;~-1[$ et sur $]2~;~+\infty[$.
\end{enumerate}

\medskip

\textbf{\large Partie B}

\medskip

On admet que la fonction $f$ de la partie $\mathbf{A}$ est définie sur $\mathbb{R}$ par 
$f(x)=\left(x^{2}- 5 x + 6\right) \mathrm{e}^{x}$.

On note $\mathcal{C}$ la courbe représentative de la fonction $f$ dans un repère.



\begin{enumerate}
\item 
	\begin{enumerate}
		\item On détermine la limite de la fonction $f$ en $+\infty$.
		
$\left .
\begin{array}{l}
\ds\lim_{x\to +\infty} \left (x^2-5x+6\right ) = \ds\lim_{x\to +\infty} x^2 = +\infty\\
\ds\lim_{x\to +\infty} \e^{x} = +\infty
\end{array}
\right \}		
\implies
\ds\lim_{x\to +\infty} \left (x^2-5x+6\right )\e^{x} = +\infty$

donc
$\ds\lim_{x\to +\infty} f(x)=+\infty$		
		
		\item On détermine la limite de la fonction $f$ en $-\infty$.
		
Pour $x\neq 0$, $f(x)=x^2\left (1-\dfrac{5}{x} +\dfrac{6}{x^2}\right )\e^{x}	
= x^2\e^{x}\left (1-\dfrac{5}{x} +\dfrac{6}{x^2}\right )$

$\left .
\begin{array}{l}
\ds\lim_{x\to -\infty} x^2 \e^{x} = 0\\
\ds\lim_{x\to -\infty} \left (1-\dfrac{5}{x} +\dfrac{6}{x^2}\right )=1
\end{array}
\right \}		
\implies
\ds\lim_{x\to -\infty} x^2\e^{x}\left (1-\dfrac{5}{x} +\dfrac{6}{x^2}\right ) = 0$

donc
$\ds\lim_{x\to -\infty} f(x)=0$
		
	\end{enumerate}
\item 
La fonction $f$ est dérivable sur $\R$ et:

$f'(x)=\left (2x-5\right )\e^{x} + \left (x^2-5x+6\right )\times \e^{x}
= \left ( 2x-5+x^2-5x+6\right )\e^{x}
= \left (x^2-3x+1\right )\e^{x}$


\item Pour déterminer le sens de variation de la fonction $f$, on étudie le signe de $f'(x)$.

Pour tout $x$, $\e^{x}>0$ donc $f'(x)$ est du signe de $x^2-3x+1$.

$\Delta=(-3)^2-4\times 1 \times 1 = 5>0$;
$x'=\dfrac{-(-3)-\sqrt{5}}{2\times 1} = \dfrac{3-\sqrt{5}}{2}$
et $x''=\dfrac{3+\sqrt{5}}{2}$

\begin{center}
{
\renewcommand{\arraystretch}{1.5}
\def\esp{\hspace*{2cm}}
$\begin{array}{|c | *{7}{c} |} 
\hline
x  & -\infty & \esp & \frac{3-\sqrt{5}}{2} & \esp & \frac{3+\sqrt{5}}{2} & \esp & +\infty \\
\hline
x^2-3x+1 &  & \pmb{+} &  \vline\hspace{-2.7pt}{0} & \pmb{-} & \vline\hspace{-2.7pt}{0} & \pmb{+} &\\
\hline
\e^x &  & \pmb{+} &  \vline\hspace{-2.7pt}{\phantom 0} & \pmb{+} & \vline\hspace{-2.7pt}{\phantom 0} & \pmb{+} &\\
\hline
f'(x) &  & \pmb{+} &  \vline\hspace{-2.7pt}{0} & \pmb{-} & \vline\hspace{-2.7pt}{0} & \pmb{+} & \\
\hline
f &  & \text{croissante} &  \vline\hspace{-2.7pt}{\phantom 0} & \text{décroissante} & \vline\hspace{-2.7pt}{\phantom 0} & \text{croissante} & \\
\hline
\end{array}$
}
\end{center}


\item La tangente $(\mathcal{T})$ à la courbe $\mathcal{C}$ au point d'abscisse 0 a pour équation réduite: \\
$y=f'(0)\left (x-0\right ) + f(0)$.

~~\textbullet~~$f(x)=\left (x^2-5x+6\right )\e^{x}$ donc $f(0)=6\e^{0}=6$

~~\textbullet~~$f'(x)=\left (x^2-3x+1\right )\e^{x}$ donc $f'(0)=1\e^{0}=1$

$\mathcal{T}$ a pour équation réduite $y=1\left (x-0\right )+6$ soit $y=x+6$.

\end{enumerate}

On admet que la fonction $f$ est deux fois dérivable sur $\mathbb{R}$. On note $f''$ la fonction dérivée seconde de la fonction $f$. On admet que, pour tout réel $x$, on a $f''(x) = (x + 1)(x - 2) \mathrm{e}^{x}$.

\begin{enumerate}[resume]
\item 
	\begin{enumerate}
		\item Pour étudier la convexité de la fonction $f$ sur $\mathbb{R}$, on étudie le signe de $f''(x)$.
		
\begin{center}
{
\renewcommand{\arraystretch}{1.5}
\def\esp{\hspace*{2cm}}
$\begin{array}{|c | *{7}{c} |} 
\hline
x  & -\infty & \esp & -1 & \esp & 2 & \esp & +\infty \\
\hline
x+1 &  & \pmb{-} &  \vline\hspace{-2.7pt}{0} & \pmb{+} & \vline\hspace{-2.7pt}{\phantom 0} & \pmb{+} &\\
\hline
x-2 &  & \pmb{-} &  \vline\hspace{-2.7pt}{\phantom 0} & \pmb{-} & \vline\hspace{-2.7pt}{0} & \pmb{+} &\\
\hline
\e^x &  & \pmb{+} &  \vline\hspace{-2.7pt}{\phantom 0} & \pmb{+} & \vline\hspace{-2.7pt}{\phantom 0} & \pmb{+} &\\
\hline
f''(x) &  & \pmb{+} &  \vline\hspace{-2.7pt}{0} & \pmb{-} & \vline\hspace{-2.7pt}{0} & \pmb{+} & \\
\hline
f &  & \text{convexe} &  \vline\hspace{-2.7pt}{\phantom 0} & \text{concave} & \vline\hspace{-2.7pt}{\phantom 0} & \text{convexe} & \\
\hline
\end{array}$
}
\end{center}
		
		\item 
Sur $[-1~;~2]$, la fonction $f$ est concave donc la courbe $\mathcal{C}$ est en dessous de ses tangentes, et en particulier en dessous de la tangente $\mathcal{T}$ car $0\in [-1~;~2]$.		
		
Donc, pour tout $x$ appartenant à l'intervalle $[-1~;~2]$, on a $f(x) \leqslant x + 6$.		
	\end{enumerate}
\end{enumerate}

\bigskip

\end{document}
